% ===================================================================
%  اللوحة نمرة ١ — المدرسة. — الثانوية. — الفصل.
%
%  Ceci n'est PAS une transcription : c'est une traduction, faite en
%  2026, en ARABE EGYPTIEN ECRIT. D'ou trois differences avec texto/io
%  et texto/fr, et trois seulement :
%
%   * pas de \nl ni de \cc — la traduction ne suit aucune ligne
%     imprimee, et les lignes de ce fichier ne sont que du confort ;
%   * pas de \begin{VUpage} — elle n'a ni feuillet ni folio ;
%   * le gras se pose sur le TERME seul, ponctuation dehors.
%
%  Les cles « %%K » sont celles de texto/io, macro pour macro pour
%  l'apparat, et les renvois \textsuperscript{(n)} visent les MEMES
%  objets de la planche, dans le MEME ordre.
%
%  DROITE A GAUCHE. Le fichier se lit dans l'ordre logique et c'est le
%  navigateur qui retourne l'affichage ; html.py pose « dir="rtl" ».
%  Les renvois restent ecrits avec leurs chiffres latins APRES le
%  substantif, comme dans les trente autres colonnes : l'algorithme
%  bidirectionnel les remet du bon cote tout seul. Ne rien inverser a
%  la main — ce qui est retourne deux fois revient a l'endroit.
%
%  POURQUOI CETTE COLONNE EXISTE, ALORS QUE « ar » EST DEJA LA. Pour
%  la meme raison que le cantonais existe a cote du chinois : si elle
%  ecrivait l'arabe standard, elle ferait double emploi a la
%  prononciation pres, et tomberait sous le motif exact qui a fait
%  ecarter le wu (texto/lingui.json, champ « ecartita »). Elle
%  s'ecrit donc en egyptien, qui a une grammaire que l'arabe standard
%  n'a pas :
%
%    ده / دي / دول      au lieu de   هذا / هذه / هؤلاء
%    مش                 au lieu de   ليس
%    ما...ش             la negation du verbe : مايعرفش
%    بـ + inaccompli    le present : بيكتب, et non يكتب
%    حـ / هـ            le futur : هيكتب, et non سيكتب
%    بتاع / بتاعة / بتوع  la possession, au lieu de la seule annexion
%    اللي               invariable, au lieu de الذي / التي / الذين
%    أوي                au lieu de   جدًا
%    كمان               au lieu de   أيضًا
%    بس                 au lieu de   لكن / فقط
%    لسه                au lieu de   ما زال
%    زي                 au lieu de   مثل
%    علشان              au lieu de   لأن / لكي
%    كام                au lieu de   كم / بضعة
%
%  ET UNE PHONOLOGIE QUI S'ECRIT : le ث devient ت ou س, le ذ devient
%  د ou ز. D'ou تلاتة et non ثلاثة, تمن et non ثمانية, تالت et non
%  ثالث, كتير et non كثير, ده et non ذاك. Ce ne sont pas des fautes
%  d'orthographe : c'est ainsi que l'egyptien s'imprime.
%
%  L'ECOLE EGYPTIENNE A SES PROPRES TITRES, ET ILS TOMBENT JUSTE SUR
%  CEUX DU LIVRET :
%
%    ناظر (71)     le proviseur — et non مدير
%    وكيل (72)     le censeur, adjoint du ناظر
%    مفتِّش (73)    l'inspecteur
%    فرّاش (77)    le concierge — celui qui porte le registre
%    الفسحة (76)   la recreation
%    شاطر          le bon eleve
%
%  Ces quatre-la ne sont pas des traductions approchees : ce sont les
%  fonctions telles qu'elles existent dans une ecole egyptienne, avec
%  la meme hierarchie que dans le lycee de 1926.
%
%  ET LE MOBILIER SCOLAIRE PARLE ITALIEN, TURC ET FRANCAIS :
%
%    ترابيزة (18)  la table    — de l'italien tavola
%    أستيكة (36)   la gomme    — de l'italien elastica
%    دواية (25)    l'encrier
%    لمبة (62)     la lampe    — et non مصباح
%    أباجورة       l'abat-jour
%    شنطة (57)     le cartable — et non حقيبة
%    بالطو         le manteau  — de « paletot »
%    مطواة (56)    le canif
%    كراسة (28)    le cahier
%    نوتة (54)     le calepin
%    فرنساوي       « francais » — et non الفرنسية
%
%  « postigo » (t01, renvoi 78) EST UN VASISTAS, ET LA PLANCHE LE DIT.
%  Verifie sur la gravure en 2026 : la verriere de la classe a deux
%  rangees de carreaux, et dans la rangee HAUTE deux petits vantaux
%  sont graves battants, gonds sur un montant vertical, ouverts vers
%  l'interieur. C'est le seul ouvrant, et il ne sert qu'a aerer. On
%  ecrit donc « شبابيك تهوية » et non un meneau ni un volet.
% ===================================================================

%%K t01-apar-0 apar
\VUsaut{2.0mm}
\VUcentre{12.6pt}{120}{كُتيِّب الشرح}
\VUsaut{3.2mm}
\VUcentre{8.4pt}{160}{بتاع}
\VUsaut{2.6mm}
\VUtitre{15.6pt}{0}{لوحات دِلماس المساعِدة}
\VUsaut{3.4mm}
\VUornamento{9pt}
\VUsaut{5.0mm}
\VUcentre{13.2pt}{90}{السلسلة الأولى}
\VUsaut{3.0mm}
\VUfilet{16mm}
\VUsaut{5.6mm}

%%K t01-apar-1 apar
\VUtitre{15.0pt}{60}{اللوحة نمرة ١}
\VUsaut{1.4mm}
\VUtitre{11.4pt}{0}{المدرسة. --- الثانوية. --- الفصل.}
\VUsaut{2.2mm}
\VUfilet{20mm}
\VUsaut{6.4mm}

%%K t01-tit-1 sub
\VUpk{10.2pt}{40}{وصف عام.}
\VUsaut{3.0mm}

%%K t01-01-1 p
1. --- اللوحة دي بتوَرِّي \VUgras{فصل} في \VUgras{ثانوية}. في الفصل ده تمن
\VUgras{ترابيزات} \textsuperscript{(18)} وتمن \VUgras{بنوك}
\textsuperscript{(32)}. على كل بنك قاعدين تلات تلامذة. و\VUgras{المدرِّس}
\textsuperscript{(7)} قاعد على \VUgras{كرسي} \textsuperscript{(8)} قدام
\VUgras{دِكّته} \textsuperscript{(6)}.

%%K t01-02-1 p
2. --- على شمال الدكة في \VUgras{دولاب} \textsuperscript{(15)} \VUgras{قزاز}.
وعلى اليمين \VUgras{السبورة} \textsuperscript{(1)} ومعاها \VUgras{الطلّاسة}
\textsuperscript{(2)} و\VUgras{الطباشير} \textsuperscript{(3)}.

%%K t01-02-2 p
وفوق الدكة متعلَّقة \VUgras{لوحات حيطة} \textsuperscript{(9, 11, 12)}
و\VUgras{خريطة} \textsuperscript{(10)}.

%%K t01-03-1 p
3. --- مش كل التلامذة قاعدين في \VUgras{أماكنهم} \textsuperscript{(17)}. يوحنا
\textsuperscript{(4)} واقف عند السبورة؛ ماسك الطلّاسة وبيمسح \VUgras{الأشكال
الهندسية}.

%%K t01-03-2 p
وبطرس جنب الدكة؛ بيلمّ \VUgras{الواجبات المكتوبة} \textsuperscript{(5)} وبيوَدِّيها
للمدرِّس.

%%K t01-04-1 p
4. --- المدرِّس \VUgras{ده} هو مدرِّس \VUgras{اللغات الحيَّة}. بيعلّم التلامذة
يتكلّموا اللغات الأجنبية ويقروها ويكتبوها \VUgras{(ألماني وإنجليزي وإسباني
وإيطالي وروسي} أو \VUgras{فرنساوي)}.

%%K t01-05-1 p
5. --- الفصل واسع أوي ومنوَّر أوي. في آخره \VUgras{شباك} عريض ليه
\VUgras{شبّاكين تهوية} \textsuperscript{(78)} و\VUgras{باب قزاز}، علشان يهوّوه
وينوّروه.

%%K t01-05-2 p
والفصل ده مش بارد. في نُصّه، علشان يدفّيه، في \VUgras{صوبة}
\textsuperscript{(46)} كويسة أوي مع \VUgras{ماسورتها} \textsuperscript{(45)}.

%%K t01-05-3 p
وفي الحاجز \VUgras{شمّاعات} \textsuperscript{(58)} متثبّتة؛ التلامذة بيعلّقوا
فيها طواقيهم وبلاطيهم.

%%K t01-tit-2 sub
\VUsaut{5.2mm}
\VUpk{10.2pt}{40}{حوش اللعب.}
\VUsaut{3.6mm}

%%K t01-06-1 p
6. --- \VUgras{الساعة} \textsuperscript{(63)} بتوَرِّي تسعة وخمس دقايق.

%%K t01-06-2 p
\VUgras{الفسحة} \textsuperscript{(76)} لسه خالصة والدرس بيبتدي تاني. اللعب
بيبقى في \VUgras{الحوش} \textsuperscript{(75)}. إحنا شايفين كام تلميذ بيلعبوا.
و\VUgras{مساعد المدرِّس} \textsuperscript{(74)} بيراقبهم.

%%K t01-07-1 p
7. --- في الحوش كمان شايفين ناس تانيين، يعني \VUgras{المفتِّش}
\textsuperscript{(73)}، و\VUgras{ناظر المدرسة} \textsuperscript{(71)}،
و\VUgras{الوكيل} \textsuperscript{(72)}.

%%K t01-07-2 p
المفتِّش بيفتِّش على المدارس، والناظر بيدير المدرسة، والوكيل بيراقب المذاكرة.

%%K t01-07-3 p
والشخص الرابع هو \VUgras{الفرّاش} \textsuperscript{(77)}. ماسك تحت باطه دفتر
علشان يكتب فيه الغايبين.

%%K t01-08-1 p
8. --- في آخر الحوش في \VUgras{أجهزة الجمباز} \textsuperscript{(70)}
و\VUgras{الورشة} بتاعة \VUgras{الشغل اليدوي} \textsuperscript{(69)}.

%%K t01-08-2 p
وعلى يمين اللوحة بنشوف \VUgras{فصول} \VUgras{الموسيقى} و\VUgras{الرسم}.

%%K t01-noto-1 noto
\VUnotes{143.2mm}{%
(*) بالنسبة \textit{لأسامي المعمودية} كمان الأحسن إنك تكتبها على الصورة
اللاتينية. دي الطريقة الوحيدة اللي تخلّصك من اختلافات ما لهاش آخر. وبعدين
تختار إزاي بين \textit{Giovanni, Jean, Johann, Jan, Hans, John, Ivan,}
وهكذا؟ هنعمل قاموس مخصوص لأسامي المعمودية؟ خلّينا ناخد من اللاتيني حالة
الرفع بتاعتها ونقول: \VUgras{Ioannes, Ioanna; Iulius, Iulia; Iakobus;}
\VUgras{Andreas; Lukas.} (Gram. Kompleta).%
}

%%K t01-tit-3 sub
\VUpk{10.2pt}{40}{وصف بالتفصيل.}
\VUsaut{2.4mm}

%%K t01-tit-4 sub
\VUpk{10.2pt}{40}{التلامذة الشاطرين.}
\VUsaut{3.0mm}

%%K t01-09-1 p
9. --- تلامذة البنك الأول على يمين الدكة هم \VUgras{هنريكوس} (*)،
و\VUgras{ستيفانوس}، و\VUgras{كارولوس}.

%%K t01-09-2 p
هنريكوس واخد باله أوي وبيحب الشغل، وبيسمع كلام المدرِّس.

على ترابيزته عنده \VUgras{كراسة} \textsuperscript{(28)}، و\VUgras{كتاب}
\textsuperscript{(29)}، و\VUgras{قاموس} \textsuperscript{(30)}، و\VUgras{مقلمة}
\textsuperscript{(31)}. كتابه فيه \VUgras{صفحات} كتير، وكل فصل فيه كذا
\VUgras{فقرة}، وكل \VUgras{سطر} فيه \VUgras{كلمات} كتير.

%%K t01-09-4 p
جاره ستيفانوس \textsuperscript{(24)} مجتهد. بيكتب في كراسته درس المرة الجاية.
\VUgras{خطّه} حلو. كتابه مقفول. إحنا شايفين \VUgras{الغلاف}
\textsuperscript{(27)} بس. وجنبه \VUgras{الفرجار} \textsuperscript{(26)}.

%%K t01-09-5 p
كارولوس \textsuperscript{(23)} ماسك كتابه بإيده، وبيفتحه علشان يراجع درسه. عنده،
زي كل تلميذ، \VUgras{دواية} \textsuperscript{(25)} قدامه. وده ولد سامع الكلام.

%%K t01-10-1 p
10. --- على الترابيزة اللي جنبها \VUgras{لودوفيكوس}، و\VUgras{أنطونيوس}،
و\VUgras{باولوس}. لودوفيكوس بيسمّع \VUgras{درسه} \textsuperscript{(22)} وبيرد على
\VUgras{أسئلة} المدرِّس. أنطونيوس \textsuperscript{(21)} مش واخد باله خالص، وساعات
المدرِّس بيعاقبه. باولوس \textsuperscript{(20)} زميل كويس، لطيف وبيحب يساعد. وهو
واخد باله أوي.

%%K t01-tit-5 sub
\VUsaut{4.2mm}
\VUpk{10.2pt}{40}{التلامذة الوحشين.}
\VUsaut{3.2mm}

%%K t01-11-1 p
11. --- ورا هنريكوس قاعد \VUgras{جورجيوس} \textsuperscript{(38)}. مش ولد وحش،
بس \VUgras{بيرغي} كتير، ومش واخد باله، وشقي. \VUgras{استخفافه}
و\VUgras{عناده} بيعملوا \VUgras{لخبطة} في الدرس، وكتير بيستاهل \VUgras{إنذار}
شديد. \VUgras{كراسة المسوّدة} \textsuperscript{(35)} بتاعته وسخة، بيلطّخها بالحبر
بـ\VUgras{القلم} \textsuperscript{(34)} بتاعه، وعلشان كده دايمًا بيستعمل
\VUgras{الأستيكة} \textsuperscript{(36)} بتاعته؛ و\VUgras{شنطة الكراريس}
\textsuperscript{(33)} بتاعته مقطّعة، علشان هو مهمل أوي. هو جايب \VUgras{علبة
الأقلام} \textsuperscript{(37)} بتاعته ليه في فصل \VUgras{اللغات الحيَّة}؟

%%K t01-11-2 p
جاره إدموندوس \textsuperscript{(39)} مش بيحبه. هو الحقيقة كسلان شوية، بس ساكت
وهادي. مش بيرغي؛ إنما بدل ما يسمع، بيفضّل يقرا \VUgras{الحواديت} اللي في
\VUgras{كتاب القراءة} بتاعه.

%%K t01-11-3 p
تالت تلميذ في البنك ده هو \VUgras{لودوفيكوس} \textsuperscript{(40)}. ده مجتهد.
بيكتب على \VUgras{ورقة} بيضا. وقدامه حطّ \VUgras{ورقة النشّاف}
\textsuperscript{(41)}.

%%K t01-12-1 p
12. --- تلامذة البنك التاني، على شمال المدرِّس، صغيّرين أوي. الأول،
\VUgras{إميليوس} الصغير، لسه مش عارف يكتب كويس، جايب \VUgras{اللوح}
\textsuperscript{(42)} بتاعه، و\VUgras{قلم الحجر} بتاعه، و\VUgras{السفنجة}
\textsuperscript{(43)} بتاعته.

%%K t01-12-2 p
وجنبه \VUgras{أوغوستوس} و\VUgras{برناردوس} \textsuperscript{(44)}. ده بيشتغل
كويس، بس \VUgras{هدومه} مهملة أوي.

%%K t01-13-1 p
13. --- وراهم في تلات \VUgras{كسالى. ألبرتوس} مش بيذاكر دروسه.
و\VUgras{ياكوبوس} \textsuperscript{(47)} بينام بدل ما يسمع. \VUgras{كسله} بيجيب له
\VUgras{عتاب} وكمان \VUgras{عقاب}. وأخيرًا \VUgras{كريستيانوس}
\textsuperscript{(48)} مش جدع خالص. \VUgras{بيتصرّف} وحش ومش بيحاول يصلّح نفسه.

%%K t01-14-1 p
14. --- جيرانه بالعكس \VUgras{مؤدَّبين}. \VUgras{إرنستوس}
\textsuperscript{(51)} بيحب النظام والترتيب، دايمًا بيستعمل \VUgras{المسطرة}
\textsuperscript{(50)} بتاعته و\VUgras{القلم الرصاص} \textsuperscript{(49)} بتاعه.
وده \VUgras{تلميذ خارجي}.

%%K t01-14-2 p
\VUgras{فرانسيسكوس} \textsuperscript{(52)} \VUgras{تلميذ داخلي}، لابس بدلة موحّدة،
وبياكل وبينام في المدرسة. وده \VUgras{زميل} كويس.

%%K t01-14-3 p
موريسيوس بيبري \VUgras{قلمه} بـ\VUgras{مطواته} \textsuperscript{(56)}، وده بيهتم
أوي.

%%K t01-14-4 p
جايب كل كتبه، و\VUgras{النحو} \textsuperscript{(55)} بتاعه، و\VUgras{النوتة}
\textsuperscript{(54)} بتاعته، وكمان \VUgras{مخدّة الكتابة} \textsuperscript{(53)}.
و\VUgras{شنطة المدرسة} \textsuperscript{(57)} بتاعته على الأرض وراه.

%%K t01-14-5 p
والترابيزتين اللي في آخر الفصل قاعد عليهم \VUgras{التلامذة الشاطرين}
\textsuperscript{(19)}.

%%K t01-tit-6 sub
\VUsaut{4.6mm}
\VUpk{10.2pt}{40}{رسم وموسيقى.}
\VUsaut{3.4mm}

%%K t01-15-1 p
15. --- في \VUgras{فصل الرسم} في \VUgras{موديلات} حلوة متعلّقة على الحيطة،
و\VUgras{لمبة} \textsuperscript{(62)} ليها \VUgras{أباجورة}، متعلّقة في السقف.
و\VUgras{المدرِّس} \textsuperscript{(60)} صبور أوي، بيصلّح \VUgras{رسومات}
\textsuperscript{(61)} التلامذة وبيعلّمهم يرسموا \VUgras{تمثال نصفي}
\textsuperscript{(59)} من الطبيعة.

%%K t01-16-1 p
16. --- \VUgras{فصل الموسيقى} شكله مشوّق أوي. فيه بيتعلّموا يقروا
\VUgras{الموسيقى}، ويسولفجوا \VUgras{النوتات} \textsuperscript{(67)} المكتوبة على
\VUgras{المدرّج} (do, re, mi, fa, sol, la, si\,=\,C. D. E. F. G. A. B.)،
ويعرفوا \VUgras{المفاتيح}، ويغنّوا \VUgras{ألحان} حلوة. والنوتة بتتحطّ على
\VUgras{الحامل} \textsuperscript{(65)}. وواحد من التلامذة \VUgras{بيلعب على
البيانو} \textsuperscript{(64)} و\VUgras{مدرِّس الموسيقى} \textsuperscript{(66)}
\VUgras{بيضبط الوزن} \textsuperscript{(68)}.

%%K t01-17-1 p
17. --- بنشوف ركن من \VUgras{الورشة} بتاعة \VUgras{الشغل اليدوي}
\textsuperscript{(69)}، اللي فيها الولاد يقدروا يتعلّموا يشتغلوا الخشب بالفارة
ويبردوا الحديد. ده مش موجود في كل الثانويات.

%%K t01-tit-7 sub
\VUsaut{5.0mm}
\VUpk{10.2pt}{40}{فصل العلوم.}
\VUsaut{3.6mm}

%%K t01-18-1 p
18. --- مدرِّس الرياضة بيعلّم التلامذة \VUgras{الهندسة والحساب والعمليات}
و\VUgras{النظام المتري}. وعلى اللوحة إحنا شايفين الأشكال الهندسية الأساسية:
\VUgras{دايرة} \textsuperscript{(a)}، و\VUgras{مربّع} \textsuperscript{(b)}،
و\VUgras{مثلث} \textsuperscript{(c)}، و\VUgras{خط} \textsuperscript{(d)}.

%%K t01-18-2 p
وشايفين كمان \VUgras{عملية}، دي مش \VUgras{جمع}، ولا \VUgras{طرح}، ولا
\VUgras{ضرب}، دي \VUgras{قسمة} \textsuperscript{(e)}.

%%K t01-19-1 p
19. --- وعلى لوحة النظام المتري إحنا شايفين: \VUgras{متر}
\textsuperscript{(a)}، و\VUgras{ميزان} \textsuperscript{(b)} و\VUgras{صنجه}،
و\VUgras{لتر} \textsuperscript{(e)}، و\VUgras{ديكالتر} \textsuperscript{(f)}،
و\VUgras{ديسيلتر}، و\VUgras{متر مكعّب} \textsuperscript{(d)}.

%%K t01-19-2 p
والتلامذة بيدرسوا كمان \VUgras{علوم الطبيعة} \textsuperscript{(11)}، وعلم الحيوان
\textsuperscript{(a)}، وعلم النبات \textsuperscript{(b)}، وعلم طبقات الأرض
\textsuperscript{(c)}، و\VUgras{التاريخ} \textsuperscript{(12)}.

%%K t01-19-3 p
وفي الدولاب أجهزة \VUgras{الطبيعة} \textsuperscript{(15)} و\VUgras{الكيميا}
\textsuperscript{(16)}.

%%K t01-19-4 p
وفوق الدولاب: \VUgras{كرة} \textsuperscript{(14)}، و\VUgras{مخروط}،
و\VUgras{كرة أرضية} \textsuperscript{(13)}.

%%K t01-19-5 p
وفوق اللوحات متعلّقة \VUgras{خريطة} بتوَرِّي أجزاء الدنيا الخمسة:
\VUgras{أوروبا} \textsuperscript{(g)}، و\VUgras{آسيا} \textsuperscript{(e)}،
و\VUgras{أفريقيا} \textsuperscript{(f)}، و\VUgras{أمريكا} \textsuperscript{(ab)}،
و\VUgras{أوقيانوسيا} \textsuperscript{(h)}، والمحيطين الكبار،
\VUgras{الباسيفيكي} \textsuperscript{(c)} و\VUgras{الأطلنطي} \textsuperscript{(d)}.

\VUsaut{6.0mm}
\VUfilet{22mm}
